\chapter{동기식 설계: 클럭, 레지스터, 조합 논리}\label{ch:sync}

릴레이식 봄베(E부)의 회로를 다시 떠올려 보자. 그 회로는 세 가지 일로 나뉘었다. 접점 나무와 이진 시프터는
제어선의 값대로 길을 이을 뿐 아무것도 기억하지 않았다. 감지 릴레이와 계수기의 슬레이브는 자기 유지로 값을
붙들었다. 그리고 값이 바뀌어도 되는 순간은 순서기가 정했다. 555가 6.3~ms마다 링 계수기를 한 칸 밀고, 그
박자에 맞춰 \code{STROBE}, \code{JUDGE}, \code{CLEAR}, \code{CK}가 나갔다(\ref{ch:relay-sequencer}장).

FPGA의 회로도 같은 세 가지로 이루어진다. 계산은 조합 논리가, 기억은 플립플롭이, 순간은 클럭이 맡는다.
다른 것은 대가를 치르는 곳이다. 릴레이식은 박자 하나를 치려고 릴레이 50개짜리 순서기 보드를 따로 만들었다.
FPGA에서는 모든 기억 소자가 처음부터 클럭 입력을 하나씩 갖고 있고, 클럭 선 하나가 칩 전체에 박자를 나눠
준다. 앞 장(\ref{ch:hdl}장)의 회로는 입력이 바뀌면 출력이 따라 바뀔 뿐 아무것도 기억하지 못했다. 이 장은
릴레이 회로에서 출발해, 디지털 회로가 기억과 시간을 다루는 표준적인 방법을 설명한다.

\section{두 종류의 부품}

\Relay 릴레이식의 회로를 ``기억이 있는가''로 나누면 둘이다. 접점 나무(\ref{sec:relay-tree}절)와 이진
시프터(\ref{sec:relay-shifter}절)는 제어선의 값에 따라 길을 열고 닫는다. 제어선이 바뀌면 릴레이가 넘어간 뒤
새 길이 생기고, 예전 길은 흔적을 남기지 않는다. 지금의 제어값만이 길을 정한다. 반면 자기 유지
회로(\ref{sec:relay-latch}절)는 자기 접점으로 자기 코일을 먹여서, 켠 신호가 사라진 뒤에도 켜진 채 남는다.
지금 켜져 있는가는 지금의 입력이 아니라 지난 일이 정한다. 테스트 레지스터의 감지 릴레이와 계수기의 슬레이브가
이 꼴이다.

디지털 회로도 똑같이 두 종류의 부품으로 이루어진다.

\begin{itemize}
\item \term{조합 논리}{combinational logic}: 현재 입력만으로 출력이 정해지는 회로. 접점 나무에 해당한다.
게이트를 엮은 것이며 기억이 없다. 입력이 바뀌면 \emph{얼마 뒤에} 출력이 바뀐다. 이 ``얼마''가 지연이다.
릴레이에서는 릴레이가 넘어가는 밀리초였고, FPGA에서는 나노초 단위다.
\item \term{레지스터}{register}: 값을 기억하는 부품. 플립플롭(flip-flop)이라는 1비트 기억 소자를 묶은
것이다. 플립플롭 하나가 자기 유지 릴레이 하나에 해당한다. 다른 점이 하나 있다. 자기 유지 릴레이는 켜는
신호가 오면 언제든 켜지고, 지우려면 유지 길을 따로 끊어야 한다. 플립플롭은 \term{클럭}{clock}이라는
주기적인 신호가 0에서 1로 올라가는 순간(상승 에지)에만 입력을 받아들여 출력으로 내보내고, 다음 상승 에지까지
그 값을 붙들고 있다. 값을 받아들이는 \emph{순간}이 부품 안에 들어 있는 것이다.
\end{itemize}

\Why 왜 값을 클럭 에지에서만 바꾸는가. 릴레이식에서 먼저 보자. 시프터의 릴레이가 넘어가는 동안 길은 믿을
수 없다. 전환 접점의 공통이 어디에도 닿지 않는 틈이 있고(표~\ref{tab:t-relay}), 릴레이마다 넘어가는 시각도
조금씩 다르다. 그래서 릴레이식은 연결이 다 자리 잡은 뒤에야 시험 전압을 걸고(\ref{sec:f-current}절), 시프터의
전환을 \code{CLEAR}의 오름부터 다음 \code{STROBE}의 오름 사이로 묶어 둔다(표~\ref{tab:t-rules}의 L9).
감지 릴레이는 그렇게 길이 멈춰 있는 \code{STROBE} 동안에만 값을 받는다.

FPGA에서도 같은 일이 나노초 규모로 일어난다. 조합 논리의 출력은 입력이 바뀐 직후에는 믿을 수 없다. 신호가
게이트를 지나는 데 시간이 걸리고 경로마다 그 시간이 달라서, 잠깐 동안 엉뚱한 값이 나타났다 사라진다. 회로
전체가 ``이제 값이 안정되었으니 받아 적자''고 합의할 기준 시각이 필요한데, 그것이 클럭 에지다. 모든 레지스터가
같은 클럭에 맞춰 값을 바꾸도록 설계하는 방식을 \term{동기식 설계}{synchronous design}라 하고, FPGA 설계는
사실상 전부 이 방식이다.

\General 그림~\ref{fig:g-dff}는 플립플롭 하나의 동작이다. 입력 \code{d}는 아무 때나 바뀌어도 되지만
출력 \code{q}는 상승 에지에서만 바뀐다. \code{d}에 잠깐 나타났다 사라진 값(두 에지 사이의 짧은 펄스)은
\code{q}에 전혀 나타나지 않는다. 감지 릴레이가 \code{STROBE} 창 밖의 일을 보지 않는 것과 같은 이치다. 다만
릴레이의 창은 밀리초 폭이고, 플립플롭의 창은 에지 앞뒤의 아주 짧은 구간으로 줄어 있다.

\begin{figure}[ht]
\centering
\begin{tikzpicture}[x=1cm,y=1cm]
  \node[reg,minimum width=14mm,minimum height=16mm] (ff) at (0,0) {};
  \node[font=\ttfamily\small,anchor=west] at ([xshift=1pt]ff.west) {d};
  \node[font=\ttfamily\small,anchor=east] at ([xshift=-1pt]ff.east) {q};
  \pic at ([xshift=0mm]ff.south) {clkmark};
  \draw[<-] (ff.west) -- ++(-0.9,0);
  \draw[->] (ff.east) -- ++(0.9,0);
  \draw[<-] (ff.south) -- ++(0,-0.5) node[below,font=\ttfamily\small] {clk};
\end{tikzpicture}
\hspace{8mm}
\begin{tikztimingtable}[timing/slope=0.1,timing/coldist=4pt,xscale=1.5,yscale=1.25,semithick]
  \code{clk} & 8{.5L .5H} \\
  \code{d}   & .2L 1.5H .5L .3H 1.2L 2H 2.3L \\
  \code{q}   & .5L 1H 2L 2H 2.5L \\
\extracode
  \begin{pgfonlayer}{background}
    \vertlines[cgray!50,dashed]{0.5,1.5,...,7.5}
  \end{pgfonlayer}
\end{tikztimingtable}
\caption{플립플롭. 왼쪽은 기호, 오른쪽은 동작이다. \code{q}는 상승 에지(점선)에서의 \code{d}를 받아 다음
에지까지 유지한다.}
\label{fig:g-dff}
\end{figure}

\section{레지스터 사이의 한 클럭}

\Relay 릴레이 계수기(\ref{sec:relay-counter}절)가 같은 문제를 이미 풀었다. 1을 더할 때 뒤집힐 비트를 정하는
올림 사슬은 비트 자신을 보고 있다. 비트가 뒤집히는 동안 사슬이 흔들리면 결과가 릴레이들의 빠르기 경쟁에
달린다. 계수기 보드는 비트마다 릴레이 둘과 두 상의 클럭으로 이를 막았다(그림~\ref{fig:r-counter}).
\code{NCK} 동안 마스터가 슬레이브의 반대값, 곧 ``뒤집으면 될 값''을 미리 만들어 두고, 슬레이브는 자기 유지로
현재 값을 붙든다. \code{CK} 동안 사슬이 닿은 비트의 슬레이브가 마스터의 값을 받는데, 사슬은 그동안 움직이지
않는 마스터의 접점으로만 이루어져 있다. \code{CK}와 \code{NCK}는 릴레이 하나의 한 극에서 나오므로 결코 동시에
켜지지 않는다(\ref{sec:relay-bbm}절). 그래서 계수기의 값은 \code{CK} 펄스에서만 움직인다.

FPGA의 플립플롭은 흔히 이 짝을 한 부품 안에 넣은 꼴로 만든다. 입력을 받아 두는 쪽과 내보내는 쪽이 클럭의
두 상에서 번갈아 열리고, 바깥에서 보면 값이 상승 에지 한 순간에 넘어간다. 릴레이식에서 보드 위에 직접 짜던
마스터--슬레이브를 FPGA의 설계자는 짤 필요가 없다. 플립플롭을 쓰면 이미 들어 있다.

이렇게 보면 동기식 회로의 모든 부분은 그림~\ref{fig:g-rcr}의 꼴로 볼 수 있다. 레지스터에서 나온 값이 조합
논리를 지나 다음 레지스터(같은 레지스터일 수도 있다)의 입력에 도착한다.

\begin{figure}[ht]
\centering
\begin{tikzpicture}[x=1cm,y=1cm]
  \node[reg,minimum height=14mm] (r1) at (0,0) {레지\\스터};
  \node[cloud,minimum width=44mm,minimum height=14mm] (c) at (4.2,0) {조합 논리\\{\footnotesize 지연 $t_{\mathrm{logic}}$}};
  \node[reg,minimum height=14mm] (r2) at (8.4,0) {레지\\스터};
  \draw[->] (r1) -- (c); \draw[->] (c) -- (r2);
  \draw[<-] (r1.south) -- ++(0,-0.6) coordinate (k1);
  \draw[<-] (r2.south) -- ++(0,-0.6) coordinate (k2);
  \draw (k1) -- (k2) node[midway,below,font=\ttfamily\small] {clk};
  \pic at (r1.south) {clkmark}; \pic at (r2.south) {clkmark};
  \draw[<->,cred] (0.6,1.15) -- node[above,note,text=cred] {이 구간의 지연이 클럭 주기보다 짧아야 한다} (7.8,1.15);
\end{tikzpicture}
\caption{동기식 회로의 기본 단위. 에지에서 출발한 값이 다음 에지 전에 도착해야 한다.}
\label{fig:g-rcr}
\end{figure}

\Why 이 꼴에서 동기식 설계의 유일한 시간 규칙이 나온다. \emph{한 에지에서 레지스터를 떠난 값은 조합
논리를 지나 다음 에지 전에 다음 레지스터에 도착해 있어야 한다.} 늦으면 다음 레지스터가 덜 계산된 값을
받아 적는다. 따라서 가장 느린 조합 논리 경로의 지연이 클럭 주기의 하한을, 곧 회로가 동작할 수 있는
최고 클럭 주파수를 정한다. 이 규칙만 지키면 조합 논리 안에서 값이 어떻게 출렁이든 상관이 없다.

릴레이식에서 같은 물음은 표~\ref{tab:t-rules}의 아홉 규칙 L1--L9가 되었다. 신호마다 릴레이 시간의 최악값을
더해 ``적어도 이만큼 기다려라''를 하나하나 따로 정해야 했다. 모두 하한이어서 기계가 느려지면 늘
안전했다(\ref{sec:t-bounds}절). 동기식 규칙도 같은 성질의 하한이다. 클럭을 늦추면 언제나 더 잘 지켜진다.
이 프로젝트의 FPGA 엔진도 처음 목표한 100~MHz에서 이 규칙을 지키지 못해 클럭을 50~MHz로
낮췄다(\ref{sec:timing}절). 다른 점은 규칙의 개수다. 동기식에서는 경로마다 같은 꼴의 규칙 하나이고, 설계
도구가 모든 경로의 지연을 계산해 확인해 준다. 설계자가 신경 쓸 시간의 단위가 ``나노초''에서 ``클럭 몇
번''으로 바뀌는 것이 동기식 설계의 이점이다.

\paragraph{클럭은 어디서 오는가.} 릴레이식은 박자를 만드는 데 보드 한 장을 썼다. 순서기의 555가 6.3~ms마다
4017을 한 칸 밀고, 열 칸이 한 위치라 한 위치가 63~ms다(\ref{sec:t-sequencer}절). 박자만은 릴레이가 아니라 전자
부품으로 만들었는데, 릴레이로 만든 지연은 최소를 보증할 수 없기 때문이다. 그리고 순서기는 그 박자 하나에서
\code{STROBE}, \code{JUDGE}, \code{CLEAR}, \code{CK}처럼 시각이 다른 신호 여럿을 따로 만들어 낸다. FPGA에서는
보드의 100~MHz 발진기에서 MMCM이라는 클럭 회로가 50~MHz를 만들고, 그 클럭 선 하나가 엔진의 모든 플립플롭에
들어간다. 한 주기는 20~ns다. 시각이 다른 신호를 클럭 선으로 따로 만들지 않는다. 클럭은 하나이고, ``이번
클럭에 무엇을 하는가''는 상태 기계가 정한다(\ref{sec:fsm}절).

\General 가장 단순한 예는 카운터다(그림~\ref{fig:g-counter}). 레지스터의 출력에 1을 더하는 조합 논리를
달고, 그 결과를 다시 레지스터의 입력에 넣는다. 에지마다 값이 1씩 커진다.

\begin{lstlisting}[language=SV,numbers=none]
logic [3:0] count;
always_ff @(posedge clk)
  if (reset) count <= 4'd0;
  else       count <= count + 4'd1;
\end{lstlisting}

\begin{figure}[ht]
\centering
\begin{tikzpicture}[x=1cm,y=1cm]
  \node[reg,minimum height=13mm,minimum width=13mm] (r) at (0,0) {\code{count}};
  \node[cloud,minimum width=16mm] (a) at (2.6,0) {$+1$};
  \draw[->] (r) -- (a);
  \draw[->] (a.east) -- ++(0.5,0) |- (1.3,1.2) -| ([xshift=-7mm]r.west) -- (r.west);
  \pic at (r.south) {clkmark};
  \draw[<-] (r.south) -- ++(0,-0.45) node[below,font=\ttfamily\small] {clk};
\end{tikzpicture}
\hspace{6mm}
\begin{tikztimingtable}[timing/slope=0.1,timing/coldist=4pt,xscale=1.9,yscale=1.25,semithick]
  \code{clk}       & 5{.5L .5H} \\
  \code{count}     & .5D{0} 1D{1} 1D{2} 1D{3} 1D{4} .5D{5} \\
  \code{count+1}   & .5D{1} 1D{2} 1D{3} 1D{4} 1D{5} .5D{6} \\
\extracode
  \begin{pgfonlayer}{background}
    \vertlines[cgray!50,dashed]{0.5,1.5,...,4.5}
  \end{pgfonlayer}
\end{tikztimingtable}
\caption{카운터. 조합 논리는 ``다음 값''을 늘 계산하고 있고, 레지스터는 에지마다 그것을 받아 ``현재 값''으로
삼는다.}
\label{fig:g-counter}
\end{figure}

카운터에서 익혀 둘 것은 이름 두 개다. 레지스터에 담긴 \emph{현재 값}과, 조합 논리가 계산해 둔 \emph{다음
값}. 동기식 회로를 읽는 것은 결국 ``다음 값이 현재 값에서 어떻게 계산되는가''를 읽는 일이다. 릴레이 계수기도
같은 꼴이다. 슬레이브가 현재 값이고, 마스터의 거울과 올림 사슬이 다음 값을 만드는 부분이며, \code{CK} 펄스가
에지의 자리에 있다.

\Ours 봄베의 기준 RTL은 이 꼴 그대로다. \code{live\_q}가 676비트의 현재 상태를 담은 레지스터이고,
\code{live\_n}이 조합 논리(코드~\ref{lst:g-spread})가 계산한 다음 상태다. 이름 끝의 \code{\_q}와
\code{\_n}은 각각 레지스터의 출력과 다음 값(next)을 뜻하는 관례다. 카운터의 ``$+1$'' 자리에
``한 걸음 퍼뜨리기''가 들어갔을 뿐이다(그림~\ref{fig:g-liveloop}).

\begin{figure}[ht]
\centering
\begin{tikzpicture}[x=1cm,y=1cm]
  \node[reg,minimum height=16mm,minimum width=16mm] (r) at (0,0) {\code{live\_q}\\{\footnotesize 676비트}};
  \node[cloud,minimum width=44mm,minimum height=16mm] (f) at (4.6,0) {한 걸음 퍼뜨리기\\{\footnotesize 스크램블러 + 대각선 보드}};
  \draw[->] (r) -- (f);
  \draw[->] (f.east) -- ++(0.6,0) |- (2.3,1.5) node[pos=0.75,above,font=\footnotesize] {\code{live\_n}} -| ([xshift=-8mm]r.west) -- (r.west);
  \pic at (r.south) {clkmark};
  \draw[<-] (r.south) -- ++(0,-0.45) node[below,font=\ttfamily\small] {clk};
\end{tikzpicture}
\caption{\code{bombe\_core}의 뼈대. 그림~\ref{fig:g-counter}의 카운터와 같은 꼴이다.}
\label{fig:g-liveloop}
\end{figure}

릴레이식과 기어식에는 \code{live\_q}에 해당하는 기억이 없다. 676가닥은 전선 그 자체이고, 전압을 걸면 이어진
선이 저절로, 사실상 즉시 같은 전압이 된다(\ref{sec:f-current}절). 기억은 테스트 케이블의 26선에 달린 감지
릴레이뿐이다. 그 퍼짐을 FPGA의 게이트로 그대로 옮기면 출력이 자기 입력으로 돌아오는 고리가 되는데, 그런 고리는
한번 1이 되면 꺼지지 않고 설계 도구가 지연을 계산할 수도 없다. 그래서 676가닥 전부에 플립플롭을 하나씩 두어
고리를 끊고, 클럭마다 한 걸음씩 퍼뜨린다(\ref{sec:f-rounds}절). 전류가 공짜로 하던 일에 676비트의 기억과
걸음 수만큼의 클럭을 치르는 셈이다(\ref{sec:f-cost}절). 이 설계 결정은 \ref{ch:rtl-core}장에서 자세히 다룬다.

\section{\texorpdfstring{\code{<=}와 \code{=}}{<=와 =}: 동시에 갱신한다는 것}

\Relay 릴레이식의 계수기들은 한 위치에 한 번, 같은 \code{CK} 펄스 안에서 함께 움직인다(\ref{sec:relay-counter}절의
``계수 규칙의 배치''). 자리들 사이에 순서가 없다. 예를 들어 $a_3$는 ``$l \ne 25$''일 때만 오르는데,
이 조건은 $l$ 자리의 마스터 접점이 만든다. $l$ 자신도 같은 펄스에서 바뀌지만, 마스터는 \code{CK} 동안 움직이지
않으므로 $a_3$가 보는 것은 언제나 펄스 직전의 $l$이다. 같은 규칙을 차례로 실행되는 Rust로 적은
코드~\ref{lst:r-counters}는 다르다. \code{l\_wraps}와 \code{m\_wraps}를 맨 앞에서 옛 값으로 계산해 두고 나서야
값들을 바꾼다. 프로그램은 옛 값을 따로 적어 두어야 하고, 릴레이 계수기는 마스터가 옛 값을 붙들고 있으므로 그럴
필요가 없다.

SystemVerilog에는 대입 기호가 둘 있다. 레지스터를 기술하는 \code{always\_ff} 안에서는 \code{<=}(논블로킹
대입)를, 조합 논리를 기술하는 \code{always\_comb} 안에서는 \code{=}(블로킹 대입)를 쓴다.

\Why 클럭 에지에서 모든 레지스터는 \emph{동시에} 값을 바꾼다. 그러므로 어떤 레지스터의 새 값을 계산할 때
쓰이는 다른 레지스터의 값은 모두 \emph{에지 직전의} 값이어야 한다. \code{<=}는 바로 그 뜻이다. ``오른쪽을
지금의 값으로 계산해 두었다가, 에지에서 한꺼번에 왼쪽에 넣는다.'' 릴레이 계수기의 마스터가 하던 일을 문법
하나로 적은 것이다. 반면 \code{=}는 조합 논리를 적을 때 ``이 전선은 이 식의 결과다''를 차례로 이어 적기 위한
것이다(코드~\ref{lst:g-scr}의 \code{x}처럼).

\General 두 레지스터의 값을 맞바꾸는 회로를 보자.

\begin{lstlisting}[language=SV,numbers=none]
always_ff @(posedge clk) begin
  a <= b;
  b <= a;
end
\end{lstlisting}

프로그램이라면 첫 줄에서 \code{a}가 덮어써지므로 임시 변수가 필요하다. 회로에서는 필요 없다. 두 줄 모두
에지 직전의 값을 읽고, 에지에서 동시에 쓴다. 그림으로는 두 레지스터의 출력이 서로의 입력에 이어져 있을
뿐이다(그림~\ref{fig:g-swap}).

\begin{figure}[ht]
\centering
\begin{tikzpicture}[x=1cm,y=1cm]
  \node[reg,minimum width=11mm] (a) at (0,0) {\code{a}};
  \node[reg,minimum width=11mm] (b) at (3,0) {\code{b}};
  \draw[->] (a.east) -- ++(0.4,0) |- (1.5,0.9) -| ([xshift=-4mm]b.west) -- (b.west);
  \draw[->] (b.east) -- ++(0.5,0) |- (1.5,-1.0) -| ([xshift=-5mm]a.west) -- (a.west);
\end{tikzpicture}
\hspace{10mm}
\begin{tikztimingtable}[timing/slope=0.1,timing/coldist=4pt,xscale=1.8,yscale=1.25,semithick]
  \code{clk} & 5{.5L .5H} \\
  \code{a}  & .5D{3} 1D{8} 1D{3} 1D{8} 1D{3} .5D{8} \\
  \code{b}   & .5D{8} 1D{3} 1D{8} 1D{3} 1D{8} .5D{3} \\
\extracode
  \begin{pgfonlayer}{background}
    \vertlines[cgray!50,dashed]{0.5,1.5,...,4.5}
  \end{pgfonlayer}
\end{tikztimingtable}
\caption{맞바꾸기. 두 레지스터가 같은 에지에서 서로의 옛 값을 받는다.}
\label{fig:g-swap}
\end{figure}

\Ours \code{bombe\_core}에서 상태를 갱신하는 부분은 다음과 같다. 세 줄의 \code{<=}는 차례로 일어나는
일이 아니라 같은 에지에서 한꺼번에 일어나는 일이다. \code{iterations + 16'd1}에 쓰인 \code{iterations}도,
조건식의 \code{live\_q}도 모두 에지 직전의 값이다. 릴레이 계수기의 자리들이 한 \code{CK} 펄스에서 함께 오르는
것과 같다.

\begin{lstlisting}[language=SV,numbers=none]
SPREAD: begin
  live_q     <= live_n;
  iterations <= iterations + 16'd1;
  if (live_n == live_q || (early_exit && test_n == FULL)) begin
    state <= IDLE;
    done  <= 1'b1;
  end
end
\end{lstlisting}

\section{``차례로''는 상태 기계로 만든다}\label{sec:fsm}

\Relay 테스트 레지스터는 한 위치마다 \code{STROBE}, \code{JUDGE}, \code{CLEAR}를 이 차례로 받아야
한다(\ref{sec:t-register}절). 붙잡고, 판정하고, 지우는 표본 유지의 순서다. 이 차례를 만드는 장치가 기계마다
있다.
\begin{itemize}
\item \textbf{릴레이식의 순서기}(\ref{sec:t-sequencer}절). 4017의 출력 Q0--Q9 가운데 늘 하나만 켜져 있고,
555의 틱마다 다음 칸으로 넘어간다. ``지금 몇 번째 단계인가''를 이 칩이 기억한다. 단계마다 무엇을 내보낼지는
Q 몇 개를 다이오드 OR로 묶은 창이 정한다. \code{STROBE}는 Q0--Q1, \code{JUDGE}는 Q3--Q4, \code{CLEAR}는
Q6--Q8이고, Q9에서 링이 주차한다(그림~\ref{fig:t-ring}). 4017은 전자 칩이지만, 그 위에 얹힌 조건은 릴레이가
붙든다. \code{CK} 결정 릴레이는 Q6에서 \code{RUN}을 보고 붙어 Q6--Q8 동안 스스로를 붙들고, 한 걸음 버튼의
\code{KSA}와 \code{KSC}는 ``한 바퀴만 돈다''를 기억한다.
\item \textbf{기어식의 슬롯 원판}(\ref{sec:t-gear}절). 상태를 기억하는 부품이 따로 없다. 주기 축이 돈 각도가
곧 ``지금 몇 번째 단계인가''이고, 세 트랙의 슬롯이 단계마다의 출력이다(그림~\ref{fig:t-slots}). 각도로 만든
상태 기계인 셈이다. 그래도 각도만으로는 알 수 없는 사실, 곧 ``이번 주기에 \code{STROBE}를 냈는가''는 따로
기억해야 해서 제어 보드는 \code{CYC}와 \code{SEEN} 두 비트를 둔다. 실제 회로에서 이 둘은 D 플립플롭
두 개다(\ref{sec:gk-gates}절).
\end{itemize}

\Why 회로는 모든 것이 동시에 동작하므로, ``먼저 이것을 하고 그다음에 저것을 한다''는 순서는 저절로
생기지 않는다. 순서가 필요하면 ``지금 몇 번째 단계인가''를 기억하는 레지스터를 두고, 단계에 따라 다른
동작을 하게 한다. 이것이 \term{상태 기계}{finite state machine, FSM}다. 4017의 켜진 출력 하나가 바로 그
기억이었다. 프로그램에서는 처리 장치가 ``지금 몇 번째 줄인가''를 대신 기억해 주지만, 회로에서는 설계자가 그
기억을 직접 만든다.

\General 버튼을 누르면 LED를 정해진 시간 동안 켜 두는 회로를 생각하자. 상태는 둘이다. 기다리는
\code{IDLE}과 켜 두는 \code{ON}. \code{IDLE}에서 버튼이 눌리면 타이머를 0으로 하고 \code{ON}으로 간다.
\code{ON}에서는 클럭마다 타이머를 올리다가 정해진 값이 되면 \code{IDLE}로 돌아간다
(그림~\ref{fig:g-fsm} 왼쪽). 상태를 담는 방법은 여럿이다. 4017처럼 상태마다 선을 하나씩 두고 하나만 켜는
방법도 있고, 상태에 번호를 매겨 이진수로 담는 방법도 있다. 어느 쪽이든 상태를 담는 것은 레지스터이고, 다음
상태를 정하는 것은 조합 논리다.

\Ours \code{bombe\_core}의 상태 기계도 상태가 둘이고 구조가 같다(그림~\ref{fig:g-fsm} 오른쪽).
\code{IDLE}에서 시작 신호가 오면 상태를 초기화하고 \code{SPREAD}로 간다. \code{SPREAD}에서는 클럭마다 한
걸음씩 퍼뜨리다가 더 변하지 않으면 완료를 알리고 \code{IDLE}로 돌아간다. 타이머 대신 ``변화가 없는가''가
끝내는 조건일 뿐이다.

릴레이식의 순서와 나란히 놓으면 차이가 보인다. \code{IDLE}에서 \code{start}를 본 에지 하나에서 세 가지가 함께
일어난다(\ref{ch:rtl-core}장). 스크램블러들이 내놓은 치환표를 레지스터에 붙잡고, 676비트를 모두 0으로 지우고,
전압을 거는 선 하나만 1로 한다. 릴레이식에서는 지우기(\code{CLEAR})와 전압 걸기(\code{STROBE})가 서로 다른
박자였다. 감지 릴레이가 다 놓일 만큼 \code{CLEAR}를 길게 끌어야 했고, 유지 버스가 다시 살아난 뒤에야 다음
\code{STROBE}를 걸 수 있었다(표~\ref{tab:t-rules}의 L5, L6). 플립플롭은 모두 같은 에지에서 바뀌므로 그런
기다림이 없다. 판정도 따로 박자를 받지 않는다. 끝난 뒤의
\code{test\_q}가 다 찼는가를 보는 조합 논리 한 줄이다. 또 하나, 릴레이식의 한 위치는 링이 한 바퀴 도는
63~ms로 늘 같지만, \code{SPREAD}는 고정점에 닿을 때까지 머물므로 위치마다 걸음 수가
다르다(\ref{sec:f-rounds}절).

\begin{figure}[ht]
\centering
\begin{tikzpicture}[x=1cm,y=1cm,
   st/.style={draw,circle,minimum size=14mm,fill=cgreenbg,draw=cgreen,font=\small\ttfamily},
   sb/.style={draw,circle,minimum size=14mm,fill=cbluebg,draw=cblue,font=\small\ttfamily}]
  \node[font=\small\bfseries] at (1.8,2.4) {일반적인 예: LED 타이머};
  \node[st] (i) at (0,0) {IDLE};
  \node[st] (o) at (3.6,0) {ON};
  \draw[->] (i) to[bend left=25] node[above,note] {버튼 눌림\\타이머 $\leftarrow 0$} (o);
  \draw[->] (o) to[bend left=25] node[below,note] {타이머 = 끝값} (i);
  \draw[->] (o) to[out=30,in=-30,looseness=5] node[right,note] {타이머 $+1$} (o);
  \begin{scope}[xshift=8.6cm]
  \node[font=\small\bfseries] at (1.8,2.4) {\code{bombe\_core}};
  \node[sb] (i2) at (0,0) {IDLE};
  \node[sb] (s2) at (3.6,0) {SPREAD};
  \draw[->] (i2) to[bend left=25] node[above,note] {\code{start}\\상태 초기화} (s2);
  \draw[->] (s2) to[bend left=25] node[below,note] {변화 없음: \code{done}} (i2);
  \draw[->] (s2) to[out=30,in=-30,looseness=5] node[right,note] {한 걸음} (s2);
  \end{scope}
\end{tikzpicture}
\caption{상태 기계 두 개. 상태를 기억하는 것도 레지스터이고, 화살표는 ``다음 상태''를 정하는 조합 논리다.}
\label{fig:g-fsm}
\end{figure}

\section{정리}

\begin{itemize}
\item 회로는 조합 논리(계산)와 레지스터(기억)로 이루어진다. 릴레이식의 접점 나무와 자기 유지 릴레이가 각각
그 짝이다.
\item 레지스터는 클럭 에지에서만 값을 바꾸고, 모든 레지스터가 동시에 바꾼다. 릴레이 계수기가
마스터--슬레이브와 \code{CK}/\code{NCK}로 만들던 성질이 플립플롭 하나에 들어 있다.
\item 설계자가 적는 것은 ``현재 값에서 다음 값이 어떻게 나오는가''다.
\item 레지스터 사이의 조합 논리는 한 클럭 주기 안에 끝나야 한다. 이 제약이 최고 속도를 정한다. 릴레이식의
L1--L9처럼 하한이어서, 클럭을 늦추면 언제나 더 잘 지켜진다.
\item 순서가 필요하면 상태 기계를 만든다. 릴레이식은 링 계수기로, 기어식은 원판의 각도로, FPGA는 상태
레지스터로 ``지금 몇 번째 단계인가''를 기억한다.
\end{itemize}

표~\ref{tab:g-sync-map}은 이 장의 대응을 한데 모은 것이다.

\begin{table}[ht]
\centering
\caption{릴레이식과 FPGA에서 계산, 기억, 시간을 맡는 것.}
\label{tab:g-sync-map}
\small
\begin{tabular}{>{\raggedright\arraybackslash}p{24mm}>{\raggedright\arraybackslash}p{58mm}>{\raggedright\arraybackslash}p{50mm}}
\toprule
& 릴레이식 & FPGA \\
\midrule
기억 없는 계산 & 접점 나무, 이진 시프터 & 조합 논리 \\
1비트 기억 & 자기 유지 릴레이(감지 릴레이, 계수기의 슬레이브) & 플립플롭 \\
값이 넘어가는 순간 & \code{CK} 펄스(마스터--슬레이브, \code{CK}/\code{NCK}) & 클럭의 상승 에지 \\
시간 기준 & 순서기의 555와 4017, 한 틱 6.3~ms & 클럭 선 하나, 50~MHz \\
시간 규칙 & L1--L9, 신호마다 따로 정한 하한 & 경로마다 ``지연 $<$ 클럭 주기'' \\
순서 & 링 계수기 Q0--Q9와 상태 릴레이 & 상태 레지스터와 상태 기계 \\
\bottomrule
\end{tabular}
\end{table}
